% ----------------------------------------------------------
\chapter{Fundamentação Teórica}\label{cap:desenvolvimento}
% ----------------------------------------------------------

\section{O Método de Monte Carlo}

O método de Monte Carlo \cite{fishman1996} é uma técnica numérica amplamente utilizada na computação para resolver problemas que envolvem incerteza, variabilidade ou integração de funções complexas por meio de simulações estocásticas. Fundamenta-se na geração de números aleatórios para modelar fenômenos que seriam difíceis ou impossíveis de tratar de forma determinística, permitindo estimar valores aproximados de integrais, probabilidades ou soluções de sistemas complexos \cite{metropolis1949}.

\section{Algoritmos Recursivos}

Os algoritmos recursivos são métodos de resolução de problemas em que uma função chama a si mesma para resolver subproblemas menores da mesma natureza. Esse tipo de algoritmo é especialmente útil quando o problema pode ser dividido em partes semelhantes, permitindo que cada chamada trate de um caso mais simples até atingir uma condição de parada ou caso base \cite{rosen2012}.

\section{Circunferência de Raio Unitário e Cálculo de \texorpdfstring{$\pi$}{pi}}

A área de uma circunferência de raio unitário no plano euclidiano é determinada pela expressão geral $A_{circulo} = \pi r^{2}$, que, para $r=1$, resulta em $A_{circulo} = \pi$.

Para calcular a área de $1/4$ da circunferência unitária usando o método de Monte Carlo, a ideia é colocar este arco dentro de um quadrado de aresta 1 e sortear pontos aleatórios. Se o ponto sorteado estiver dentro de $1/4$ da circunferência, marca-se um acerto.

A probabilidade $P$ exata de um ponto cair na circunferência é igual à razão entre as áreas, conforme a \autoref{eq:Eq_probabilidade}:
\begin{equation}\label{eq:Eq_probabilidade}
    P = \frac{A_{circulo}}{A_{quadrado}} = \frac{\pi/4}{A_{quadrado}}.
\end{equation}

Como a área do quadrado é unitária ($A_{quadrado}=1$), tem-se pela \autoref{eq:Eq_pi_exato}:
\begin{equation}\label{eq:Eq_pi_exato}
    \pi = 4 P A_{quadrado} = 4P.
\end{equation}

Assim, a partir de $N$ pontos sorteados, obtém-se uma estimativa iterativa para $\pi$ dada pela \autoref{eq:Eq_pi_aprox}:
\begin{equation}\label{eq:Eq_pi_aprox}
    \pi_{k} \approx 4 \times \frac{N_{acertos}^{(k)}}{N}.
\end{equation}

\section{Geradores de Números Pseudoaleatórios}

A precisão destes métodos depende dos geradores de números aleatórios. Em Python, a biblioteca padrão disponibiliza o módulo \texttt{random}, cujo mecanismo é baseado no algoritmo \textit{Mersenne Twister}. Os valores são classificados como pseudoaleatórios pois resultam de processos determinísticos dependentes de uma semente inicial.

\section{Erros Absoluto, Relativo e Relativo Percentual}

Em métodos numéricos utilizam-se três principais tipos de erro para avaliar a precisão:
\begin{alineas}
    \item \textbf{Erro absoluto:} é a diferença entre o valor aproximado ($x_{aprox}$) e seu valor verdadeiro ($x_{verd}$), conforme a \autoref{eq:Eq_erro_abs}:
    \begin{equation}\label{eq:Eq_erro_abs}
        \epsilon_{a} = |x_{aprox} - x_{verd}|.
    \end{equation}
    
    \item \textbf{Erro relativo:} relaciona o erro absoluto com o valor verdadeiro, mostrando a proporção do desvio, dado pela \autoref{eq:Eq_erro_rel}:
    \begin{equation}\label{eq:Eq_erro_rel}
        E_{r} = \frac{|x_{aprox} - x_{verd}|}{|x_{verd}|}.
    \end{equation}
    
    \item \textbf{Erro relativo percentual:} é o erro relativo expresso em porcentagem, obtido pela \autoref{eq:Eq_erro_perc}:
    \begin{equation}\label{eq:Eq_erro_perc}
        \epsilon_{r\%} = E_{r} \times 100\%.
    \end{equation}
\end{alineas}